% HD-R10. GPT-6 Astra Pro /root, inherited effort; 21 November 2026.
\section{HD-R10: the general sixfold secant trace obstruction}
\label{sec:general-sixfold-secant-trace}

Let $X$ be a complex abelian sixfold, let $H$ and $B$ be real
$(1,1)$ classes, and suppose that $H$ is nondegenerate of Hermitian
signature $(3,3)$. Put $s=ic$, where $c>0$. Consider an actual
perfect complex $E$ with
\[
 \operatorname{Ext}^{<0}(E,E)=0,\qquad
 \operatorname{ch}(E)=e^B
       \bigl(\alpha e^{sH}+\bar\alpha e^{-sH}\bigr),
                         \quad\alpha\ne0.                 \tag{R10.1}
\]
Let $G\subset X\times\widehat X$ be any faithful finite point
subgroup with an actual coherent translation/tensor action and
a specified equivariant structure on $E$. The trivial group is
included. Write $\mathcal C=D_{\mathrm{perf}}(X)^G$.

At least one of the two actual categorical trace maps
\[
 \sigma_E^0:\operatorname{End}_{\mathcal C}(E)\longrightarrow HH_0(\mathcal C),
 \qquad
 \sigma_E^2:\operatorname{Ext}^2_{\mathcal C}(E,E)
                                      \longrightarrow HH_{-2}(\mathcal C)
                                                               \tag{R10.2}
\]
has a nonzero kernel. Moreover every weakly $G$-semiregular member
has $\dim\operatorname{End}_{\mathcal C}(E)\ge3$. In particular
every equivariantly simple complex (R10.1) fails weak
$G$-semiregularity. The assertion covers
arbitrary ranks, cohomological amplitudes and the stated coherent
finite point actions. It supplies the exact remaining degree-zero
trace gate for a nonsimple weakly semiregular object of this form;
it does not assert that every such nonsimple object is excluded.

\paragraph{1. The actual Hochschild ring and the two graph maps.}
Put $T=H^{0,1}(X)$ and
\[
 U=T\oplus H^{1,0}(X)^*,\qquad HH^*(X)=\bigwedge^*U.
                                                               \tag{R10.3}
\]
The group-difference isomorphism $(x,y)\mapsto(x,y-x)$ carries
the diagonal to $X\times\{0\}$. Its formal Koszul resolution,
tensored with $H^*(\mathcal O_X)$, gives precisely this exterior
Yoneda algebra. Evaluation on the actual object is its graded
ring homomorphism
\[
 \operatorname{ev}_E:HH^*(X)\longrightarrow\operatorname{Ext}^*(E,E).
                                                               \tag{R10.4}
\]
The characteristic homomorphism and the trace identity
$\sigma_E^k\operatorname{ev}_E^k(a)=a\cdot\operatorname{ch}(E)$
are Buchweitz--Flenner, Adv. Math.217 (2008),243--281,
Proposition6.2.1, Proposition6.3.1, Corollary6.3.2 and
Theorem6.4.2; author version arXiv:math/0606730, pp.29--33.
The multiplicative abelian identification and all normalization
maps were rechecked in HD-N19, paragraphs1--3.

Multiplication by $e^{-B}$ conjugates the action of
$u=(\eta,v)\in U$ to the action with antiholomorphic component
$\eta+\iota_vB$. We use this invertible linear change of the
actual Hochschild generators to remove $B$ in the calculations.
It extends to the exterior algebra in(R10.3). Thus every operator
below remains evaluation of an actual Hochschild class. A real
$B$ is not replaced by a line bundle or by an asserted functor.

After this coordinate change the two maps are
\[
 \pi_\pm:U\longrightarrow T,\qquad
                \pi_\pm(\eta,v)=\eta\pm s\iota_vH.
                                                               \tag{R10.5}
\]
The map $(\pi_+,\pi_-)$ is invertible, since $s\ne0$ and
$H^\sharp:H^{1,0*}\to T$ is invertible. Let
$U_+=\ker\pi_-$ and $U_-=\ker\pi_+$; these are complementary
six-dimensional subspaces. Their restrictions to the other
$\pi$ are isomorphisms. Acting on $e^{\pm sH}$ gives exactly
\[
       u_1\cdots u_k e^{\pm sH}
        =\pi_\pm(u_1)\cdots\pi_\pm(u_k)e^{\pm sH}.
                                                               \tag{R10.6}
\]
Indeed holomorphic contraction kills each antiholomorphic factor
already produced, and contraction of the exponential supplies
the next factor. In particular every mixed product
$U_+\wedge U_-$ kills the entire character.

For every $1\le k<6$, the images
$\bigwedge^kT\,e^{sH}$ and $\bigwedge^kT\,e^{-sH}$ have
zero intersection. Their lowest components identify the two
putative representatives as a common $\xi\in H^{0,k}$;
their next components give $2sH\wedge\xi=0$.
Write $H=\sum_{i=1}^6z_i b_i$, with both sets bases.
The coefficient at $z_i$ is $b_i\wedge\xi$. For $k<6$, every
nonzero coefficient of $\xi$ has a missing basis index $i$,
so not all these coefficients can vanish. Therefore
\[
 \operatorname{rank}c_E^k=2\binom6k\quad(1\le k<6),
 \qquad (\operatorname{rank}c_E^1,
         \operatorname{rank}c_E^2,
         \operatorname{rank}c_E^3,
         \operatorname{rank}c_E^4)=(12,30,40,30),
                                                               \tag{R10.7}
\]
where $c_E^k=\sigma_E^k\operatorname{ev}_E^k$.

\paragraph{2. Degree-two injection determines the degree-four image.}
The exact finite-action comparison in HD-N12, paragraphs1--3,
identifies the forgetful Hochschild map with an isomorphism
$HH_*(\mathcal C)\to HH_*(X)$. Nonidentity translation graphs
miss the diagonal; the nontrivial flat-tensor sectors have zero
cohomology. These are the actual sector and induction/forgetful
maps, and their proofs have no dimension-four restriction.
Every point functor fixes $HH^*(X)$, so the natural evaluation
(R10.4) takes values in invariant self-Ext under every coherent
linearization. Consequently weak $G$-semiregularity is precisely
injectivity of the ambient trace on invariant degree-two Ext.

Assume this degree-two trace is injective. The mixed evaluations
have zero trace by(R10.6), and are invariant. They are therefore
zero as actual Ext classes. Multiplicativity of(R10.4) makes
evaluation factor through the explicit graded algebra
\[
 R=\frac{\bigwedge(U_+\oplus U_-)}{(U_+\wedge U_-)},
 \qquad R_k=\bigwedge^kU_+\oplus\bigwedge^kU_-
                                      \quad(k\ge1).
                                                               \tag{R10.8}
\]
The last equality follows by sorting each exterior monomial:
every monomial using both summands lies in the indicated ideal,
and the two pure exterior algebras project back from the quotient.
Thus the actual degree-four evaluation has image at most30.
Equation(R10.7) gives the reverse bound, so its image has
dimension exactly30.

Fix an invariant holomorphic volume form on $X$. The
Calabi--Yau sixfold Serre pairing restricts perfectly to
invariants, since averaging is an exact projection and all point
functors preserve this volume. It identifies the dual of
$\sigma_E^2$ with $\operatorname{ev}_E^4$ after the volume-form
Hochschild identification. The literal pairing is
$\operatorname{tr}(\xi\operatorname{ev}_E^4(a))$ for
$\xi\in\operatorname{Ext}^2(E,E)^G$; this is the trace/evaluation
adjunction used in HD-N11 and HD-N19. Injectivity of
$\sigma_E^2$ makes this adjoint evaluation surjective. Writing
$e_k=\dim\operatorname{Ext}^k(E,E)^G$, we obtain
\[
               e_2=e_4=30,\qquad e_1\ge12,\qquad e_3\ge40.
                                                               \tag{R10.9}
\]
No degree-four or degree-three trace injection is assumed.

\paragraph{3. The simple-object Euler contradiction.}
Put $V_H=\int_XH^6/6!$. The product of the Hermitian eigenvalues
has sign $(-1)^3$, so $V_H<0$. Todd class is one, and the $B$
factors cancel in the dual product. The only top-degree terms
in the product of the two secant characters are
$|\alpha|^2(e^{2sH}+e^{-2sH})$. Hence
\[
 \chi(E,E)=128|\alpha|^2s^6V_H>0,
 \qquad
       \chi_{\mathcal C}(E,E)=\frac{\chi(E,E)}{|G|}>0.
                                                               \tag{R10.10}
\]
The second equality is HD-N12's exact eigen-Euler calculation:
it treats pure tensor sectors by their connected torsion-line
cover and translation sectors by finite-etale descent of the
actual endomorphism complex. It applies to coherent scalar
twists as proved there and does not assert ordinary descent of
$E$. In the first equality $s^6=-c^6<0$.

For an equivariantly simple $E$, one has $e_0=e_6=1$.
Gluability and Serre duality leave only degrees zero through six.
Equations(R10.9)--(R10.10) would give
\[
 0<\chi_{\mathcal C}(E,E)
     =2e_0-2e_1+2e_2-e_3
                 \le2-24+60-40=-2.                         \tag{R10.11}
\]
This contradiction proves the simple-object assertion for every
stated coherent finite point action.
Without simplicity the same calculation gives
$0<\chi_{\mathcal C}(E,E)\le2e_0-4$, so $e_0\ge3$.

\paragraph{4. The exact mixed-operator Lie algebra for summand characters.}
We now prove the stronger two-trace assertion(R10.2).
The following explicit Clifford calculation supplies the needed
constraint on the Chern character of each actual summand.
Let $v$ be any real even cohomology class, of Hodge degree zero,
annihilated by every mixed operator $U_+\wedge U_-$ after the
$B$-shear. Its conjugate mixed operators annihilate it as well,
because $e^{-B}v$ is real.

Work in the full complex Clifford space
\[
 W=\bigl(H^1(X,\mathbb R)\oplus H^1(X,\mathbb R)^*\bigr)_{\mathbb C},
 \qquad c(a+\xi)=a\wedge+\iota_\xi.
\]
Use the pairing $Q(a+\xi,b+\eta)=\xi(b)+\eta(a)$, so that
$c(w)c(w')+c(w')c(w)=Q(w,w')$.
Let $A$ and $A'$ be the annihilators of $e^{sH}$ and
$e^{-sH}$ respectively. They are the actual transverse
isotropic graphs $a=\mp s\iota_\xi H$, each of dimension12.
The Clifford pairing identifies $A'$ with $A^*$.
Complex conjugation interchanges $A$ and $A'$. In the notation
of paragraph1 put
\[
 P=U\cap A=U_-,\quad R=U\cap A'=U_+,
 \qquad Q=\bar U\cap A=\bar R,\quad S=\bar U\cap A'=\bar P.
                                                               \tag{R10.12}
\]
Thus $A=P\oplus Q$ and $A'=R\oplus S$, with all four
summands of dimension6. The pairing vanishes on $P\times R$
and on $Q\times S$, since $U$ and $\bar U$ are isotropic.
Its other two blocks are nondegenerate, identifying
$R=Q^*$ and $S=P^*$.

Under this literal pairing,
$[c(p)c(r),c(a)]=Q(r,a)c(p)$ for $a\in A$, since
$Q(p,a)=Q(p,r)=0$. Thus the operators $c(p)c(r)$ for
$p\in P,r\in R$ act on $A$ as the exact off-diagonal
maps $a\mapsto Q(r,a)p$ in $\operatorname{Hom}(Q,P)$.
Their conjugates $c(s)c(q)$ act as
$a\mapsto-Q(s,a)q$ in $\operatorname{Hom}(P,Q)$.
Taking the negatives of the latter gives the opposite matrix
units with coefficient one. To check generation explicitly,
choose bases of $P,Q$ and their paired duals. These operators
are the matrix units $E_{i\mu}$ and $E_{\mu i}$.
Their brackets give $E_{ij}$, $E_{\mu\nu}$ and the diagonal
differences $E_{ii}-E_{\mu\mu}$. These span all trace-zero
matrices. Therefore their Lie algebra is exactly
$\mathfrak{sl}(A)$, through the stated Clifford representation.

The relative Fock map is an actual vector-space isomorphism
\[
 \bigwedge^*A'\longrightarrow H^*(X,\mathbb C),\qquad
 a'_1\cdots a'_k\longmapsto
       c(a'_1)\cdots c(a'_k)e^{sH}.                         \tag{R10.13}
\]
After multiplying by $e^{-sH}$, its degree-$k$ leading term is
the exterior power of the invertible map $2sH^\sharp$.
All other terms have lower exterior degree, which proves
invertibility. This is the same relative Fock morphism
explicitly constructed in HD-R08, paragraph3.
The $\mathfrak{sl}(A)$-invariants in this model are precisely
the degree-zero and degree12 axes. Indeed, in every intermediate
exterior degree, a basis monomial has an index $i$ which occurs
and an index $j$ which does not; the trace-zero diagonal matrix
$E_{ii}-E_{jj}$ gives it nonzero weight. All diagonal matrices
act diagonally, so their common zero-weight space in those
degrees is zero. The two extreme axes are fixed by every
trace-zero matrix.

The first axis in(R10.13) is $e^{sH}$ and the second is the
one-dimensional space annihilated by $A'$, namely
$\mathbb C e^{-sH}$. Consequently the original class has
the exact form
\[
            v=e^B\bigl(\beta e^{sH}+\bar\beta e^{-sH}\bigr).
                                                               \tag{R10.14}
\]
Reality gives the displayed conjugate coefficients. This
calculation identifies the morphisms and generated Lie algebra;
it does not infer(R10.14) from an abstract group name.

\paragraph{5. Simultaneous trace injection forces a single summand.}
Suppose both maps in(R10.2) are injective. The exact
Atiyah generalized-trace and coherent kernel-filtration
argument of HD-R04, paragraph2, proves that $\sigma_E^0$
kills commutators and nilpotents. Its primary additivity input
is Carmeli--Cnossen--Ramzi--Yanovski, Forum Math. Sigma13
(2025),e93, Theorem5.27 and Remark5.25, pp.57--58.
The point functors preserve the coherent-sheaf truncations
and the induced filtrations, as required in that argument.
It follows that
\[
 \operatorname{End}_{\mathcal C}(E)=\mathbb C^r,\qquad
                      E=\bigoplus_{j=1}^r E_j,              \tag{R10.15}
\]
with actual equivariant idempotent summands and
$\operatorname{End}_{\mathcal C}(E_j)=\mathbb C$.
The trace of each nonzero identity idempotent is nonzero by
degree-zero injectivity. Under the exact forgetful Hochschild
isomorphism it becomes the ordinary nonzero character
$v_j=\operatorname{ch}(E_j)$.

Paragraph2 makes every mixed degree-two evaluation on $E$
zero. Composing with the actual inclusion and projection of
$E_j$ gives zero evaluation on $E_j$. Taking its trace shows
that $v_j$ is annihilated by all these mixed operators.
It is a real rational Hodge class, so paragraph4 applies.
Thus every $v_j$ has(R10.14), with $\beta_j\ne0$.
Equation(R10.7), applied to this character, gives rank30
for its degree-two trace--evaluation composite. Its
evaluation is invariant under the given equivariant
structure. The actual diagonal inclusions of derived Hom
complexes give a direct sum of these invariant self-Ext
spaces inside $\operatorname{Ext}^2_{\mathcal C}(E,E)$.
Consequently
\[
                         30=e_2\ge30r.                      \tag{R10.16}
\]
Here $r\ge1$, since the whole character is nonzero. Therefore
$r=1$, and(R10.15) says $E$ is equivariantly simple.
Paragraph3 gives the contradiction. This proves(R10.2).

\paragraph{6. Scope of the reduced construction gate.}
HD-R08 already forces the imaginary-secant form(R10.1) for
every nondegenerate single-class polynomial source whose whole
functorial orbit reaches a sixfold Weil character. HD-R10 now
excludes every gluable equivariantly simple realization, including
all ranks and all the specified faithful coherent point actions.
For any remaining gluable weakly semiregular realization of this
form, the degree-zero equivariant trace has a nonzero kernel.
This conclusion is stronger than the ordinary simple rank-two
restriction of HD-N19 and its finite-action numerical constraints.
It retains arbitrary nonsimple complexes with that exact trace
gate as unresolved. The character calculation gives no effective
detecting cycle for a new determinant component.

All maps used here are the actual diagonal characteristic map,
equivariant forgetful/induction maps, Serre adjunction and
idempotent-summand maps on the given perfect complex. The relative
Fock map is a cohomological comparison, not an asserted algebraic
correspondence. In the programme's framework this is a restriction
on the perfect-complex/K-theory supplier needed before the bounded
cycle family in condition A; it introduces no positivity-to-cycle
step. The algebraic characteristic and pushforward operations retain
the SGA6 and GRR interpretation already recorded in the framework
section, while the real Clifford calculation is stated only in its
proved Betti/Hodge realization.

\paragraph{7. The same exact obstruction for every pure Weil sixfold character.}
There is a direct target statement, requiring no polynomial source
or Fourier transport. Let $(X,K,\theta)$ be any polarized Weil
sixfold of signature $(3,3)$, with any imaginary quadratic field
and any determinant. Every actual gluable perfect complex with
whole character
\[
                \operatorname{ch}(E)=e^Bt,\qquad
                       0\ne t\in W_K(X)_{\mathbb Q},          \tag{R10.17}
\]
has the same two-trace obstruction(R10.2) for every specified
faithful coherent finite point action. Every weakly semiregular
member has invariant endomorphism dimension at least three.

Here are the actual exterior maps proving the statement.
Choose holomorphic eigenbases $a_1,a_2,a_3$ and $b_1,b_2,b_3$
for the two field embeddings, orthonormal for the positive
polarization form. Put $y_i=\bar a_i$ and $z_i=\bar b_i$.
The eigenspaces are orthogonal because the Rosati involution
on $K$ is conjugation. With this normalization,
\[
 \theta=\frac{i}{2}\sum_{j=1}^3(a_jy_j+b_jz_j),\qquad
 \tau_+=a_1a_2a_3z_1z_2z_3,\quad
 \tau_-=b_1b_2b_3y_1y_2y_3.
                                                               \tag{R10.18}
\]
Conjugation gives $\overline{\tau_+}=-\tau_-$, and hence
$t=\lambda\tau_+-\bar\lambda\tau_-$ with $\lambda\ne0$.
The actual Hochschild generator spaces are
\[
 U_-=\langle\iota_{a_j},y_j\wedge\rangle,\qquad
 U_+=\langle\iota_{b_j},z_j\wedge\rangle.                    \tag{R10.19}
\]
Each space kills one Weil volume and acts freely by the indicated
contractions and wedges on the other. In degree $k$, the
$\binom6k$ choices give distinct exterior monomials. The two
sets have no common monomial for $k<6$: a common monomial would
have neither an $a_j$ nor a $b_j$, forcing all three holomorphic
contractions on each side, and would then require all three
additional antiholomorphic wedges, forcing $k=6$.
This proves(R10.7) directly, and the mixed block kills $t$.
The $B$-shear is exactly the linear change in paragraph1.

The full annihilator of $\tau_+$ is
$\langle a_j,z_j,\iota_{b_j},\iota_{y_j}\rangle$; that of
$\tau_-$ is its transverse paired complement. Their
intersections with $U$ and $\bar U$ are precisely the four
six-dimensional blocks of(R10.12), with the two signs relabelled.
The same literal Clifford products are the two off-diagonal
matrix blocks. Their commutators generate $\mathfrak{sl}_{12}$,
and wedging and contracting on these displayed bases gives
the relative Fock isomorphism with invariant axes
$\mathbb C\tau_+$ and $\mathbb C\tau_-$. Thus the summand
argument of paragraphs4--5 applies to these exact Weil axes.

Finally, reordering the twelve one-forms in(R10.18) gives
\[
 \tau_+\tau_-=64\frac{\theta^6}{6!},\qquad
 t^2=-128|\lambda|^2\frac{\theta^6}{6!}.                    \tag{R10.20}
\]
For example the permutation from the order in $\tau_+\tau_-$
to $a_1y_1a_2y_2a_3y_3b_1z_1b_2z_2b_3z_3$ has27 inversions;
the factor $(i/2)^6=-1/64$ supplies the other sign.
Derived duality multiplies degree-six cohomology by $(-1)^3$,
and the $B$ factors cancel. Therefore
\[
             \chi(E,E)=-\int_Xt^2
                  =128|\lambda|^2\int_X\theta^6/6!>0.      \tag{R10.21}
\]
The exact eigen-Euler map divides this by $|G|$.
Equations(R10.8)--(R10.11) and(R10.15)--(R10.16) now prove
all stated conclusions for(R10.17). This calculation retains
every field and determinant and identifies the actual maps;
no cohomological comparison is promoted to an unconstructed
geometric equivalence. The non-curve target
$e^B(P(\theta)+\lambda t)$ retains its separate polynomial
terms: paragraph7 concerns exactly $P=0$.
